\section{Problem setting and data}\label{sec:problem}

\subsection{Operational context}\label{sec:context}

The application that motivates this work is berth planning at the Port of San
Antonio, Chile. The port comprises five terminals, DP World, STI, PANUL, QC and
EPSA, each of which plans its own berth schedule independently but all of which
share a single port-entry navigation channel. A schedule assigns each incoming
vessel to a compatible berth and to a start time, with the aim of reducing vessel
waiting and improving berth utilization. This is the berth allocation problem
(BAP), for which the operations research literature offers a mature taxonomy of
discrete and continuous formulations \citep{bierwirth2010survey}. The load-bearing
input to any such schedule is the expected service time: how long each vessel will
occupy its berth. A planner cannot decide when the next vessel may enter a berth
without an estimate of when the current occupant will leave, so the quality of the
schedule is bounded by the quality of that estimate.

Berths are heterogeneous. Some accept only particular cargoes, most plainly the
liquid-bulk berth, which serves liquid-bulk vessels alone, and others are
constrained by vessel draft and length. Rather than encode these restrictions from
a berth specification sheet, we derive the vessel-to-berth compatibility matrix
from the data, treating a terminal and vessel-type pair as compatible when it
co-occurs in the historical record. A subset of calls that the port designates
\emph{service} vessels, its own term for committed or priority calls, carry a hard
no-wait reservation: they are entitled to berth on arrival. Enforcing that
guarantee as a hard constraint can render a weekly instance infeasible whenever a
predicted service time pushes a reserved berthing window past its arrival. Because
our training procedure must solve the scheduling model repeatedly under many
different predicted service times, we relax the guarantee during training to a soft
tardiness penalty on reserved vessels, which keeps every solve feasible while
preserving the priority these calls are owed. Planning is carried out one calendar
week at a time; the optimizer receives a fixed seven-day set of calls and does not
roll its horizon forward as new information arrives.

The service time a schedule depends on is one segment of a longer port call, and
naming that segment precisely matters for what the predictor is asked to learn. A
call proceeds through arrival at the anchorage, pilot boarding, mooring, cargo
operations, unmooring and departure. The two early stages have very different
character. Waiting at anchorage can last days, especially for dry-bulk vessels,
whose time at anchor has a median of about 207~hours in our records; this waiting
is largely a consequence of congestion and berth availability rather than a
property of the vessel. Pilot boarding, by contrast, marks the actual entry to the
port and is followed by a transit to the first mooring that is
consistent across calls, about 42~minutes. The berth-occupation stage that begins at mooring is
what the schedule commits berth capacity to, and it is what we predict. Prediction
of this quantity matters precisely because it is the input on which the rest of the
plan is built, and because its errors do not stay local. An under-estimate leads
the optimizer to admit the following vessel too early, and the resulting overlap
propagates down the berth queue; an over-estimate merely leaves a berth idle for a
while. We develop this \emph{cascade asymmetry} and its consequences for the
training objective in Section~\ref{sec:method}; here it is enough to note that
service-time errors are both consequential and asymmetric in their effect.

\subsection{Data}\label{sec:data}

Our data come from a vessel-calls workbook maintained by the port, covering
December~2019 through August~2025. The raw file holds 5{,}605 calls, which we clean
to 5{,}597 usable records by removing entries with inconsistent or missing
timestamps. Each record describes one call: the terminal and berth used, the
vessel and its attributes, the ports of origin and destination, the commercial
service, and the timestamps that bound each stage of the call. The port records
sixteen distinct vessel types, which we collapse into seven groups for modelling:
Container, Dry Bulk, Vehicle Carrier, Liquid Bulk, General Cargo, Passenger and a
residual Other. Container calls are the most frequent and occupy a berth for about
36~hours on average; dry-bulk calls are the longest at roughly 74~hours; vehicle
carriers average about 36~hours and liquid-bulk calls about 22~hours. These group
averages differ enough to make vessel type an informative feature, but they also
conceal wide within-group variation that a point estimate must confront.

The prediction target is the berth-occupation time, denoted $\tau$ and measured in
hours. We define it as the interval between the first mooring and the last
unmooring of a call, so that it captures the full time the vessel holds the berth,
including cargo operations and any idle time between them, and excludes both the
anchorage wait that precedes entry and the channel transit that precedes mooring.
The empirical distribution of $\tau$ is strongly right-skewed: most calls are
short, a long tail of calls is very long, and this shape informs both the choice of
error measure and the model families we consider in Section~\ref{sec:method}.

The predictor may use only information available before the vessel berths, since a
service-time estimate that relied on post-berthing data could not be produced when
the schedule is built. Our feature table has seventeen such features. Several are
recorded directly: the berth or site (\emph{Sitio}), the vessel-type group, a
COVID-era indicator marking the period of disrupted operations, the gross tonnage
(TRG), the arrival draft, and a draft difference. The timing of the berthing
decision enters through cyclical encodings, the sine and cosine of the berthing
hour, day of week, and month, which let the model represent daily, weekly and
seasonal patterns without imposing an artificial discontinuity at the wrap-around.
The remaining features are high-cardinality categoricals, namely the origin and
destination ports, the service route, the shipping line, and the agency, which we
encode by target encoding rather than one-hot expansion.

Because several features summarize a vessel's own history or its group's recent
behaviour, feature engineering is where target leakage is most likely to enter, and
we construct these features causally. Per-vessel historical statistics are computed
over an expanding window shifted by one visit, so that a given row sees only that
vessel's earlier calls and never its current or future ones. Cross-vessel group
averages are computed under the same discipline: for a call being scheduled, a
group average includes only same-group calls whose service had already been
completed before the current vessel's berthing decision, never calls still in
progress or yet to arrive. Timestamps that would reveal the label are held out of
the training table entirely. Together these rules aim to reproduce, at training
time, exactly the information a planner would hold at the moment the schedule is
built.

\subsection{Limitations of the data}\label{sec:datalimits}

Three features of the data deserve to be stated plainly, because a careful reader
should weigh them as threats to the validity of what follows. First, the
draft-difference feature is computed as arrival draft minus departure draft, and the
departure draft is known only after the vessel has finished loading or discharging,
that is, after the service whose duration we are trying to predict. This feature is
therefore a suspected source of target leakage, and to the extent the models lean on
it the reported predictive accuracy may be optimistic. We flag this explicitly as an
open issue requiring verification and a leakage-free re-fit that either removes the
feature or replaces it with a pre-berthing proxy. Second, the exact
feature-engineered training table used to fit the models is not fully reproducible
from the released pipeline scripts, so an independent researcher could reconstruct
the procedure but not necessarily the identical inputs. Third, the underlying
vessel-call data are proprietary to the port and cannot be released publicly, which
limits external replication to the code and the aggregate statistics reported here.
None of these is disqualifying, but each bounds how strongly the empirical results
should be read, and we return to the leakage concern in particular when interpreting
model performance.

Finally, we note a gap in the literature that this study begins to address. Machine
learning models for vessel service and turnaround times have been studied in other
port systems \citep{yan2024predicting, chu2025vessel}, but we are not aware of prior
machine learning work on service-time prediction for Chilean ports specifically. The
present study is, to our knowledge, a first step in that direction.
